\documentclass[11pt,a4paper]{article}
\usepackage[T1]{fontenc}
\usepackage[utf8]{inputenc}
\usepackage{amsmath,amssymb,amsthm}
\usepackage[margin=1in]{geometry}
\usepackage[colorlinks=true,linkcolor=blue,urlcolor=blue]{hyperref}
\theoremstyle{plain}
\newtheorem{theorem}{Theorem}
\newtheorem{lemma}[theorem]{Lemma}
\newtheorem{proposition}[theorem]{Proposition}
\newtheorem{corollary}[theorem]{Corollary}
\theoremstyle{definition}
\newtheorem{remark}[theorem]{Remark}

\title{A quadrisection identity between OEIS A084703 and OEIS A089928}
\author{Adrian Perez Fontelles\\ \small Independent researcher}
\date{31 August 2026}
\begin{document}
\maketitle

\begin{abstract}
OEIS A084703 carries an empirical observation of A.~Ratushnyak that its $n$-th term
equals the $(4n-2)$-nd term of OEIS A089928. It is true, and the proof is a single line
once both entries' own closed forms are put side by side: the exponent $4n$ produced by
the quadrisection turns $(1\pm\sqrt2)^{4n}$ into $(17\pm12\sqrt2)^{n}$, and the parity
term $(-1)^{\lfloor (4n-2)/2\rfloor}$ is constantly $-1$. No generating function is
needed, and in particular the generating function recorded on A089928 is not used --- as
noted below, it does not reproduce that entry's own terms.
\end{abstract}

\noindent\small 2020 Mathematics Subject Classification. 11B37, 11B39, 11D09.\normalsize

\section{The two sequences and the conjecture}

OEIS A084703 is ``Squares k such that 2*k+1 is also a square.''. It has offset $0$ and begins
\[
0, 4, 144, 4900, 166464, 5654884, 192099600,\ \dots
\]
OEIS A089928 is ``a(n) = 2*a(n-1) + 2*a(n-3) + a(n-4), with a(0)=1, a(1)=2, a(3)=4, a(4)=10.''. It has offset $0$ and begins
\[
1, 2, 4, 10, 25, 60, 144, 348,\ \dots
\]

The entry A084703 carries the following comment:
\begin{quote}\small
Empirical: a(n) = A089928(4*n-2), for n > 0. --- \emph{Alex Ratushnyak}, Apr 12 2013
\end{quote}
As of the ``Last modified'' line on the live entry (Jan 05 2026 20:36:39, revision 68) the
statement is still recorded as unproved, and nothing on either entry records it as settled
either way.

Written out, the statement to be proved is
\begin{equation}\label{eq:conj}
a(n)\;=\;b(4n-2)\qquad\text{for all } n\ge 1,
\end{equation}
where $a$ is A084703 and $b$ is A089928.

\section{What each entry states as fact}

A084703 states, not as a conjecture,
\begin{quote}\small
a(n) = ((17+12*sqrt(2))\^{}n + (17-12*sqrt(2))\^{}n - 2)/8.
\end{quote}
that is,
\begin{equation}\label{eq:a}
a(n)\;=\;\frac{\bigl(17+12\sqrt2\bigr)^{n}+\bigl(17-12\sqrt2\bigr)^{n}-2}{8} .
\end{equation}
A089928 states, not as a conjecture,
\begin{quote}\small
a(n) = ( (1+sqrt(2))\^{}(n+2) + (1-sqrt(2))\^{}(n+2) + 2*(-1)\^{}floor(n/2) )/8.
\end{quote}
that is,
\begin{equation}\label{eq:b}
b(m)\;=\;\frac{\bigl(1+\sqrt2\bigr)^{m+2}+\bigl(1-\sqrt2\bigr)^{m+2}
+2(-1)^{\lfloor m/2\rfloor}}{8} .
\end{equation}
Each was checked against its own entry's DATA at every published index before use:
\eqref{eq:a} against all $18$ terms of A084703 and \eqref{eq:b} against all
$30$ terms of A089928, in exact arithmetic in $\mathbb{Z}[\sqrt2]$. These two
lines are the only input taken from the entries.

\begin{remark}\label{rem:gf}
A089928 also carries the line ``G.f.: 1/((1+2*x)*(1-2*x-x\^{}2))''. That is not the
generating function of A089928: the stated denominator expands to $1-5x^{2}-2x^{3}$,
whose reciprocal begins $1,0,5,\dots$, whereas the entry begins $1,2,4,\dots$. The
factor $(1+2x)$ should read $(1+x^{2})$, since
$(1-2x-x^{2})(1+x^{2})=1-2x-2x^{3}-x^{4}$, matching the recurrence in the entry's
name. Nothing below uses that line; it is recorded only so that a reader checking the
entry is not misled.
\end{remark}

\section{The proof}

\begin{lemma}\label{lem:pow}
$\bigl(1+\sqrt2\bigr)^{4}=17+12\sqrt2$ and $\bigl(1-\sqrt2\bigr)^{4}=17-12\sqrt2$.
\end{lemma}

\begin{proof}
$\bigl(1\pm\sqrt2\bigr)^{2}=1\pm2\sqrt2+2=3\pm2\sqrt2$, and
$\bigl(3\pm2\sqrt2\bigr)^{2}=9\pm12\sqrt2+8=17\pm12\sqrt2$.
\end{proof}

\begin{lemma}\label{lem:sign}
For every integer $n\ge1$, $(-1)^{\lfloor (4n-2)/2\rfloor}=-1$.
\end{lemma}

\begin{proof}
$4n-2$ is even, so $\lfloor (4n-2)/2\rfloor=2n-1$, which is odd.
\end{proof}

\begin{theorem}
For every integer $n\ge1$, $a(n)=b(4n-2)$, where $a$ is OEIS A084703 and $b$ is OEIS
A089928. This is the statement of Section~1.
\end{theorem}

\begin{proof}
Put $m=4n-2$ in \eqref{eq:b}. Then $m+2=4n$, so by Lemma~\ref{lem:sign}
\[
b(4n-2)=\frac{\bigl(1+\sqrt2\bigr)^{4n}+\bigl(1-\sqrt2\bigr)^{4n}-2}{8} .
\]
By Lemma~\ref{lem:pow}, $\bigl(1\pm\sqrt2\bigr)^{4n}
=\Bigl(\bigl(1\pm\sqrt2\bigr)^{4}\Bigr)^{n}=\bigl(17\pm12\sqrt2\bigr)^{n}$, so
\[
b(4n-2)=\frac{\bigl(17+12\sqrt2\bigr)^{n}+\bigl(17-12\sqrt2\bigr)^{n}-2}{8},
\]
which is exactly $a(n)$ by \eqref{eq:a}.
\end{proof}

\begin{remark}
The identity in fact holds at $n=0$ as well: both sides are $0$, since $b(-2)$ is not
defined by the entry but \eqref{eq:b} evaluated at $m=-2$ gives $(1+1-2)/8=0=a(0)$. The
restriction $n>0$ in the comment is therefore only a statement about which indices of
A089928 are published.
\end{remark}

\section{Verification}

Three checks, each independent of the algebra above.

First, formula \eqref{eq:a} was evaluated in exact arithmetic in $\mathbb{Z}[\sqrt2]$ and
compared against every one of the $18$ published terms of A084703; the values agree
exactly. A formula that reproduced only some of them would not have been used.

Second, formula \eqref{eq:b} was evaluated the same way and compared against every one of
the $30$ published terms of A089928; the values agree exactly. It was also checked
against the recurrence $b(m)=2b(m-1)+2b(m-3)+b(m-4)$ named on that entry.

Third, the identity \eqref{eq:conj} was checked directly on the published terms, without
either formula: for every $n$ with $1\le n\le 7$ --- the range for which both $a(n)$ and
$b(4n-2)$ are published --- the two integers agree.

\begin{thebibliography}{9}
\bibitem{oeis} The OEIS Foundation, \emph{The On-Line Encyclopedia of Integer Sequences},
\url{https://oeis.org}, 2026.
\bibitem{hw} G.~H.~Hardy and E.~M.~Wright, \emph{An Introduction to the Theory of
Numbers}, 6th ed., Oxford University Press, 2008.
\bibitem{stanley} R.~P.~Stanley, \emph{Enumerative Combinatorics}, Volume 1, 2nd ed.,
Cambridge University Press, 2011.
\end{thebibliography}

\end{document}
